\section{应记应背}

高考语文注重积累，本部分涵盖古诗文默写 64 篇、文言文 120 实词、18 虚词及文化常识，是卷面基础分的保障。

\subsection{古诗文默写必背 64 篇}

高中 14 篇 + 初中 50 篇，涵盖新课标推荐背诵篇目。以下按类别分组：

\subsubsection{高中必背篇目（14 篇）}

\begin{itemize}
  \item \textbf{劝学}（荀子）—— 学不可以已；青，取之于蓝而青于蓝
  \item \textbf{逍遥游}（庄子）—— 北冥有鱼，其名为鲲；鲲之大，不知其几千里也
  \item \textbf{师说}（韩愈）—— 师者，所以传道受业解惑也
  \item \textbf{阿房宫赋}（杜牧）—— 六王毕，四海一；蜀山兀，阿房出
  \item \textbf{赤壁赋}（苏轼）—— 壬戌之秋，七月既望
  \item \textbf{氓}（《诗经》）—— 氓之蚩蚩，抱布贸丝
  \item \textbf{离骚}（屈原）—— 长太息以掩涕兮，哀民生之多艰
  \item \textbf{蜀道难}（李白）—— 噫吁嚱，危乎高哉！蜀道之难，难于上青天
  \item \textbf{登高}（杜甫）—— 风急天高猿啸哀，渚清沙白鸟飞回
  \item \textbf{锦瑟}（李商隐）—— 锦瑟无端五十弦，一弦一柱思华年
  \item \textbf{琵琶行}（白居易）—— 浔阳江头夜送客，枫叶荻花秋瑟瑟
  \item \textbf{虞美人}（李煜）—— 春花秋月何时了，往事知多少
  \item \textbf{念奴娇·赤壁怀古}（苏轼）—— 大江东去，浪淘尽，千古风流人物
  \item \textbf{永遇乐·京口北固亭怀古}（辛弃疾）—— 千古江山，英雄无觅孙仲谋处
\end{itemize}

\subsubsection{初中必背篇目（50 篇 选列）}

\begin{itemize}
  \item \textbf{论语}十二则 —— 学而时习之，不亦说乎
  \item \textbf{鱼我所欲也}（孟子）—— 鱼，我所欲也；熊掌，亦我所欲也
  \item \textbf{生于忧患，死于安乐}（孟子）—— 故天将降大任于是人也
  \item \textbf{曹刿论战}（《左传》）—— 肉食者鄙，未能远谋
  \item \textbf{邹忌讽齐王纳谏}（《战国策》）—— 吾妻之美我者，私我也
  \item \textbf{出师表}（诸葛亮）—— 先帝创业未半而中道崩殂
  \item \textbf{桃花源记}（陶渊明）—— 晋太元中，武陵人捕鱼为业
  \item \textbf{三峡}（郦道元）—— 自三峡七百里中，两岸连山，略无阙处
  \item \textbf{马说}（韩愈）—— 世有伯乐，然后有千里马
  \item \textbf{陋室铭}（刘禹锡）—— 山不在高，有仙则名
  \item \textbf{小石潭记}（柳宗元）—— 从小丘西行百二十步
  \item \textbf{岳阳楼记}（范仲淹）—— 先天下之忧而忧，后天下之乐而乐
  \item \textbf{醉翁亭记}（欧阳修）—— 醉翁之意不在酒，在乎山水之间也
  \item \textbf{爱莲说}（周敦颐）—— 出淤泥而不染，濯清涟而不妖
  \item \textbf{记承天寺夜游}（苏轼）—— 庭下如积水空明
  \item \textbf{送东阳马生序}（宋濂）—— 余幼时即嗜学
  \item \textbf{关雎}（《诗经》）—— 关关雎鸠，在河之洲
  \item \textbf{蒹葭}（《诗经》）—— 蒹葭苍苍，白露为霜
  \item \textbf{观沧海}（曹操）—— 东临碣石，以观沧海
  \item \textbf{饮酒}（陶渊明）—— 结庐在人境，而无车马喧
\end{itemize}

\subsubsection{默写易错字辨析}

\begin{center}
\begin{tabular}{c|c|c}
\textbf{篇目} & \textbf{易错字} & \textbf{正确写法} \\ \hline
《劝学》 & 輮以为轮 & 注意「輮」不写作「柔」 \\
《逍遥游》 & 抟扶摇而上 & 「抟」不写作「博」 \\
《赤壁赋》 & 壬戌 & 「戌」不写作「戍」 \\
《阿房宫赋》 & 廊腰缦回 & 「缦」不写作「漫」 \\
《琵琶行》 & 幽咽泉流 & 「咽」不写作「烟」 \\
《离骚》 & 方圜 & 「圜」通「圆」 \\
\end{tabular}
\end{center}

\newpage

\subsection{文言文 120 实词}

以下为高考高频 120 实词中的重点词汇（按考频排序）：

\subsubsection{高频实词（前 30）}

\begin{center}
\begin{tabular}{c|c|c|c}
\textbf{实词} & \textbf{本义} & \textbf{常见义项} & \textbf{例句} \\ \hline
爱 & 喜爱 & 吝惜 & 齐国虽褊小，吾何爱一牛 \\
安 & 安定 & 怎么、哪里 & 燕雀安知鸿鹄之志哉 \\
被 & 被子 & 通「披」，穿 & 将军身被坚执锐 \\
倍 & 加倍 & 通「背」，违背 & 倍道而行 \\
本 & 树根 & 本来、推究 & 人穷则反本 \\
比 & 并列 & 等到 & 比至南郡 \\
兵 & 兵器 & 军队、战争 & 兵不血刃 \\
病 & 疾病 & 困苦不堪 & 向吾不为斯役，则久已病矣 \\
察 & 观察 & 明察、了解 & 明足以察秋毫之末 \\
朝 & 早晨 & 朝廷、朝见 & 朝服衣冠 \\
曾 & 竟 & 通「增」，增加 & 曾不能毁山之一毛 \\
乘 & 驾车 & 趁着、量词 & 因利乘便 \\
诚 & 诚实 & 确实、如果 & 诚如是，则霸业可成 \\
除 & 台阶 & 授予官职 & 除臣洗马 \\
辞 & 言辞 & 推辞、告别 & 辞不就职 \\
从 & 跟随 & 听从、由 & 沛公旦日从百余骑来见项王 \\
殆 & 危险 & 大概、几乎 & 殆有神护者 \\
当 & 相抵 & 应当、正对 & 木兰当户织 \\
道 & 道路 & 道理、说 & 何可胜道也哉 \\
得 & 得到 & 能够、应该 & 意气扬扬，甚自得也 \\
度 & 量长短 & 估计、法度 & 试使山东之国与陈涉度长絜大 \\
非 & 不对 & 不是、除非 & 登高而招，臂非加长也 \\
复 & 回来 & 再、恢复 & 师道之不复可知矣 \\
负 & 背、担负 & 凭借、辜负 & 秦贪，负其强 \\
盖 & 遮盖 & 大概、发语词 & 盖失强援，不能独完 \\
故 & 原因 & 所以、旧 & 既克，公问其故 \\
固 & 坚固 & 本来、一定 & 臣固知王之不忍也 \\
顾 & 回头看 & 只是、反而 & 顾吾念之 \\
归 & 女子出嫁 & 返回、归附 & 众士慕仰，若水之归海 \\
国 & 国家 & 国都、地区 & 去国怀乡 \\
\end{tabular}
\end{center}

\subsubsection{一词多义辨析}

\begin{itemize}
  \item \textbf{举}：举起（举手长劳劳）$\mid$ 全（举世无双）$\mid$ 推举（众议举宠为督）$\mid$ 攻占（戍卒叫，函谷举）
  \item \textbf{绝}：断（绝子孙）$\mid$ 极（绝妙好辞）$\mid$ 横渡（而绝江河）$\mid$ 隔绝（率妻子邑人来此绝境）
  \item \textbf{属}：连接（属引凄异）$\mid$ 撰写（属予作文以记之）$\mid$ 隶属（十三学得琵琶成，名属教坊第一部）$\mid$ 类（有良田美池桑竹之属）
  \item \textbf{数}：数目（数口之家）$\mid$ 计算（不可遍数）$\mid$ 屡次（数见不鲜）$\mid$ 命运（则胜负之数）
  \item \textbf{间}：中间（其间千二百里）$\mid$ 间隔（间以诗记所遭）$\mid$ 秘密地（间以诗记所遭）$\mid$ 参与（肉食者谋之，又何间焉）
\end{itemize}

\subsection{文言文 18 虚词}

\begin{center}
\begin{tabular}{c|c|c|c}
\textbf{虚词} & \textbf{词性} & \textbf{常见用法} & \textbf{难点提示} \\ \hline
而 & 连词 & 表转折／并列／递进／修饰 & 注意与「以」的区别 \\
何 & 副/代 & 怎么、什么、为什么 & 疑问句中常作宾语前置 \\
乎 & 语气/介 & 吗、于 & 形容词词尾（浩浩乎） \\
乃 & 副/代 & 于是、竟然、才、你的 & 判断句中表「是」 \\
其 & 代/副 & 他/它、难道、大概 & 祈使句中表「一定」 \\
且 & 副/连 & 并且、将近、况且 & 且……且……表并列 \\
若 & 动/代/连 & 如、你、如果 & 比得上（不若） \\
所 & 助/名 & ……的地方 & 为……所（被动） \\
为 & 介/动 & 因为、替、做、被 & 介词读 wèi，动词读 wéi \\
焉 & 代/助 & 哪里、于此、……的样子 & 句末语气助词 \\
也 & 助 & 判断句标志、停顿、呢 & 不译或译「啊」 \\
以 & 介/连 & 用、因为、按照、来 & 以……为（把……当作） \\
因 & 介/副 & 凭借、于是、因为 & 因势利导 \\
于 & 介 & 在、到、从、比、被 & 常见固定结构：于是 \\
与 & 介/连/助 & 和、给予、赞同 & 通「欤」表疑问 \\
则 & 连/副 & 就、那么、却 & 则……则……表对比 \\
者 & 助 & ……的人/事/物 & 定语后置标志 \\
之 & 代/助/动 & 他、的、到、主谓之间取消独立性 & 宾语前置标志 \\
\end{tabular}
\end{center}

\newpage

\subsection{文化常识}

\subsubsection{官职变动类}

\begin{itemize}
  \item \textbf{授官}：除、拜、授、征、辟、举、起
  \item \textbf{升官}：擢、迁（右迁）、陟、加、拔
  \item \textbf{降官}：谪、贬、左迁、放、出（出官）
  \item \textbf{免官}：罢、黜、夺、免
  \item \textbf{兼职}：兼、领、摄、行、署
  \item \textbf{调动}：迁、转、调、徙
\end{itemize}

\subsubsection{科举制度}

\begin{itemize}
  \item \textbf{院试}（童试）$\rightarrow$ 生员（秀才）
  \item \textbf{乡试}（秋闱）$\rightarrow$ 举人（第一名称「解元」）
  \item \textbf{会试}（春闱）$\rightarrow$ 贡士（第一名称「会元」）
  \item \textbf{殿试} $\rightarrow$ 进士（一甲三名：状元、榜眼、探花）
\end{itemize}

\subsubsection{天文历法}

\begin{itemize}
  \item \textbf{干支纪年}：十天干 + 十二地支，60 年一循环
  \item \textbf{月相}：朔（初一）$\rightarrow$ 望（十五）$\rightarrow$ 晦（三十/廿九）
  \item \textbf{四时}：春夏秋冬，孟仲季各分一月
  \item \textbf{时辰}：一昼夜 12 时辰，对应 24 小时
\end{itemize}

\subsubsection{宗法礼仪}

\begin{itemize}
  \item \textbf{宗庙}：天子七庙、诸侯五庙、大夫三庙
  \item \textbf{丧礼}：斩衰、齐衰、大功、小功、缌麻（五服）
  \item \textbf{年龄称谓}：总角（童年）$\rightarrow$ 及笄（女子 15）$\rightarrow$ 弱冠（男子 20）$\rightarrow$ 花甲（60）$\rightarrow$ 古稀（70）$\rightarrow$ 耄耋（80--90）
\end{itemize}

\subsubsection{典籍常识}

\begin{itemize}
  \item \textbf{四书}：《大学》《中庸》《论语》《孟子》
  \item \textbf{五经}：《诗经》《尚书》《礼记》《周易》《春秋》
  \item \textbf{二十四史}：从《史记》到《明史》
  \item \textbf{史书体例}：编年体（《左传》）$\mid$ 纪传体（《史记》）$\mid$ 国别体（《国语》）
  \item \textbf{经史子集}：四部分类法
\end{itemize}
